\documentclass[10pt]{article}
\usepackage[paperwidth=384pt,paperheight=900pt,margin=10pt]{geometry}
\usepackage[T1]{fontenc}
\usepackage{lmodern}
\usepackage[utf8]{inputenc}
\usepackage{amsmath,amssymb}
\usepackage{array}
\usepackage{enumitem}
\usepackage{xcolor}
\usepackage{microtype}

\setlength{\parindent}{0pt}
\setlength{\parskip}{4pt}
\setlist[itemize]{leftmargin=12pt}
\setlist[enumerate]{leftmargin=14pt}

\sloppy
\allowdisplaybreaks
\emergencystretch=2em

\begin{document}

\newcommand{\RepeatedQuestion}[1]{%
\par\medskip
\noindent\textbf{Original question repeated for this part}\par
#1
\medskip
}

\newcommand{\uppgiftOneQuestion}{%
Låt $A$ och $B$ vara två händelser i ett och samma utfallsrum.
Sannolikheten för att ingen av händelserna $A$ eller $B$ inträffar
är $60\,\%$, medan sannolikheten för att både $A$ och $B$ inträffar
är $20\,\%$. Dessutom är sannolikheten $10\,\%$ för att händelsen
$B$ inträffar utan att händelsen $A$ inträffar.

\begin{enumerate}[label=\alph*)]
\item Är $A$ och $B$ oberoende händelser? \hfill (4 p)
\item Beräkna den betingade sannolikheten för händelsen $B$ givet
att händelsen $A$ inte inträffar. \hfill (2 p)
\end{enumerate}
}

\newcommand{\uppgiftTwoQuestion}{%
Låt följande funktion vara given.

\[
f(x)=
\begin{cases}
ax^2(2-x),
& \text{om } 0\leq x\leq 2,\\
0,
& \text{annars.}
\end{cases}
\]

\begin{enumerate}[label=\alph*)]
\item Bestäm talet $a$ sådant att $f$ är en täthetsfunktion.
\hfill (1,5 p)
\end{enumerate}

Givet att talet $a$ är sådant att $f$ är en täthetsfunktion,
låt $X$ vara en stokastisk variabel vars fördelning beskrivs av $f$.

\begin{enumerate}[label=\alph*),resume]
\item Bestäm fördelningsfunktionen för $X$.
\hfill (1,5 p)
\item Beräkna $P(X>1)$.
\hfill (1 p)
\item Beräkna väntevärdet och standardavvikelsen för $X$.
\hfill (2 p)
\end{enumerate}

\textit{Kommentar:}
Om du inte har löst a), så kan du istället använda
täthetsfunktionen
\[
f(x)=\frac{3x^2(2-x)}{4}
\]
för $0\leq x\leq 2$ för att lösa b)--d).
}

\newcommand{\uppgiftThreeQuestion}{%
Låt $X \in \operatorname{Exp}(0,04)$ och
$Y \in \operatorname{Exp}(0,05)$ vara oberoende stokastiska variabler
och sätt
\[
Z=\min\{X,Y\}.
\]

\begin{enumerate}[label=\alph*)]
\item Bestäm fördelningen för $Z$.
\hfill (4 p)
\item Beräkna $P(Z\geq 10)$.
\hfill (1 p)
\item Bestäm väntevärdet för $Z$.
\hfill (1 p)
\end{enumerate}
}

\newcommand{\uppgiftFourQuestion}{%
\begin{enumerate}[label=\alph*)]
\item
I en förening genomförs en årlig pengainsamling, där medlemmar
ombeds att bidra med endera 50 kr eller 100 kr. Under de föregående
åren har $46\,\%$ valt att inte lämna något bidrag och andelen bidrag
om 50 kr har varit dubbelt så stor som andelen bidrag om 100 kr.
Låt $X$ vara lika med det bidraget från en slumpmässigt vald medlem
baserat på föregående år. Bestäm fördelningen för $X$ och beräkna
dess väntevärde samt standardavvikelse.
\hfill (2 p)

\item
Ett taxibolag prissätter taxiresor på följande sätt: varje resa har
en startavgift på 57 kr och sedan tillkommer 14 kr per km samt
9 kr per minut under resan. Enligt en undersökning är bolagets
innerstadsresor $X$ km långa och varar i $Y$ minuter, där
$X\in N(3,1)$ och $Y\in N(8,2)$. Beräkna sannolikheten för att en
slumpmässigt vald innerstadsresa kostar mellan 150 kr och 200 kr.
\hfill (4 p)
\end{enumerate}
}

\newcommand{\uppgiftFiveQuestion}{%
I ett slumpmässigt försök räknar en basketspelare hur många lyckade
straffkast hen kan göra i rad innan det första misslyckade
straffkastet. Låt $p$ vara sannolikheten för ett lyckat straffkast
och antag att alla straffkast är oberoende av varandra. Antalet
lyckade straffkast i försöket är då en stokastisk variabel $X$ med
följande sannolikhetsfunktion.

\[
p_X(k)=p^k(1-p),
\qquad
k=0,1,2,3,\ldots
\]

\begin{enumerate}[label=\alph*)]
\item
Låt $x_1,x_2,\ldots,x_n$ vara ett observerat stickprov från
fördelningen för $X$. Härled ML-skattningen
$p^*_{\mathrm{obs}}$ av $p$.
\hfill (5 p)

\item
Basketspelaren genomförde försöket vid 8 separata tillfällen med
följande resultat: $2,0,9,1,6,7,5,3$. Beräkna ML-skattningen av
$p$ utifrån detta resultat.
\hfill (1 p)
\end{enumerate}
}

\newcommand{\uppgiftSixQuestion}{%
Ett nytt parkeringshus har byggts i Grönköping med hopp om att locka
fler personer till butikerna i centrum. Kostnaden för parkeringshuset
ska på sikt täckas av intäkterna från parkeringsavgifter. Under
30 veckodagar noterades de dagliga intäkterna från parkeringsavgifter,
vilka anses utgöra ett observerat stickprov från en normalfördelning.
Stickprovets medelvärde och standardavvikelse är 1260 kr respektive
150 kr.

\begin{enumerate}[label=\alph*)]
\item
Bestäm ett $90\,\%$ konfidensintervall för den genomsnittliga dagliga
intäkten från parkeringsavgifter baserat på stickprovet ovan.
\hfill (3 p)
\end{enumerate}

Den konsult som var rådgivande när bygget av parkeringshuset
planerades förutsåg att den genomsnittliga dagliga intäkten från
parkeringsavgifter kommer att vara 1280 kr.

\begin{enumerate}[label=\alph*),resume]
\item
Ange lämpliga hypoteser $H_0$ och $H_1$ för att pröva konsultens
hypotes. Motivera ditt svar!
\hfill (1 p)

\item
Genomför hypotesprövning enligt b) med felrisk $10\,\%$ och baserat
på det observerade stickprovet ovan. Verkar konsultens hypotes rimlig
eller kan den förkastas?
\hfill (2 p)
\end{enumerate}
}

\section*{Uppgift 1}

\subsection*{Original question}
Låt $A$ och $B$ vara två händelser i ett och samma utfallsrum.
Sannolikheten för att ingen av händelserna $A$ eller $B$ inträffar
är $60\,\%$, medan sannolikheten för att både $A$ och $B$ inträffar
är $20\,\%$. Dessutom är sannolikheten $10\,\%$ för att händelsen
$B$ inträffar utan att händelsen $A$ inträffar.

\begin{enumerate}[label=\alph*)]
\item Är $A$ och $B$ oberoende händelser? \hfill (4 p)
\item Beräkna den betingade sannolikheten för händelsen $B$ givet
att händelsen $A$ inte inträffar. \hfill (2 p)
\end{enumerate}

\subsection*{Compiled notes from Yacoub + Obid}

\subsubsection*{1. Identify the Given Information}
First, we translate the problem statements into
mathematical notation:

\begin{itemize}
\item The probability that neither $A$ nor $B$ occurs is
$60\,\%$:
\[
P(A^c \cap B^c)=0.60.
\]

\item This implies the probability that at least one occurs is:
\begin{align*}
P(A \cup B)
&= 1-0.60 \\[2mm]
&= 0.40.
\end{align*}

\item The probability that both $A$ and $B$ occur is
$20\,\%$:
\[
P(A \cap B)=0.20.
\]

\item The probability that $B$ occurs without $A$ occurring is
$10\,\%$:
\[
P(B \cap A^c)=0.10.
\]
\end{itemize}

\subsubsection*{2. Calculate Individual Probabilities}
To answer the specific questions, we need to find
$P(A)$ and $P(B)$.

\textbf{Find $P(B)$:}

The event $B$ consists of two disjoint parts:
($B$ occurring with $A$) and ($B$ occurring without $A$).

\begin{align*}
P(B)
&= P(A \cap B)
   + P(B \cap A^c)
\end{align*}

\begin{align*}
P(B)
&= 0.20+0.10 \\[2mm]
&= 0.30
\end{align*}

\textbf{Find $P(A)$:}

Using the Addition Rule:

\begin{align*}
P(A \cup B)
&= P(A)+P(B)-P(A \cap B).
\end{align*}

\begin{align*}
0.40
&= P(A)+0.30-0.20 \\[2mm]
0.40
&= P(A)+0.10 \\[2mm]
P(A)
&= 0.30
\end{align*}

\RepeatedQuestion{\uppgiftOneQuestion}
\subsubsection*{a) Are $A$ and $B$ independent events?}
\textbf{Reasoning:}
Two events are independent if and only if

\[
P(A \cap B)=P(A)\cdot P(B).
\]

\begin{enumerate}
\item We know
\[
P(A \cap B)=0.20.
\]

\item Calculate the product of their individual probabilities:
\begin{align*}
P(A)\cdot P(B)
&= 0.30 \cdot 0.30 \\[2mm]
&= 0.09
\end{align*}

\item Compare the results:
\[
0.20 \neq 0.09
\]
\end{enumerate}

\textbf{Answer:}
No, $A$ and $B$ are \textbf{not independent} because the
probability of them occurring together ($20\,\%$) is not equal
to the product of their individual probabilities ($9\,\%$).

\RepeatedQuestion{\uppgiftOneQuestion}
\subsubsection*{b) Calculate the conditional probability for event $B$ given that event $A$ does not occur.}
\textbf{Reasoning:}
We are looking for $P(B\mid A^c)$.
The formula for conditional probability is:

\begin{align*}
P(B\mid A^c)
&= \frac{P(B \cap A^c)}{P(A^c)}
\end{align*}

\begin{enumerate}
\item Find $P(A^c)$:
\begin{align*}
P(A^c)
&= 1-P(A) \\[2mm]
&= 1-0.30 \\[2mm]
&= 0.70
\end{align*}

\item Apply the formula:

Using $P(B \cap A^c)=0.10$ from the given info:

\begin{align*}
P(B\mid A^c)
&= \frac{0.10}{0.70} \\[2mm]
P(B\mid A^c)
&= \frac{1}{7} \\[2mm]
&\approx 0.143
\end{align*}
\end{enumerate}

\textbf{Answer:}
The conditional probability is $1/7$
(or approximately $14.3\,\%$).

\subsubsection*{Extra expanded solution from added notes}

\textbf{Translate the information}

Neither $A$ nor $B$ happens:

\[
P(A^c\cap B^c)=60\,\%.
\]

Therefore the chance that at least one happens is:

\begin{align*}
P(A\cup B)
&=100\,\%-60\,\% \\[2mm]
&=40\,\%.
\end{align*}

Both $A$ and $B$ happen:

\[
P(A\cap B)=20\,\%.
\]

$B$ happens but not $A$:

\[
P(B\cap A^c)=10\,\%.
\]

\textbf{Find $P(B)$}

$B$ can happen in two ways:

\begin{itemize}
\item $B$ happens and $A$ does not happen,
\item $B$ happens together with $A$.
\end{itemize}

So:

\begin{align*}
P(B)
&=
P(B\cap A^c)
+P(A\cap B) \\[2mm]
&=
10\,\%+20\,\% \\[2mm]
&=30\,\%.
\end{align*}

\textbf{Find $P(A)$}

Use the formula from the notes:

\[
P(A\cup B)=P(A)+P(B\cap A^c).
\]

Plug in:

\begin{align*}
40\,\%
&=
P(A)+10\,\% \\[2mm]
P(A)
&=
30\,\%.
\end{align*}

\RepeatedQuestion{\uppgiftOneQuestion}
\textbf{a) Check independence}

For independence:

\[
P(A\cap B)=P(A)\cdot P(B).
\]

Check with the values:

\begin{align*}
0.20
&\stackrel{?}{=}
0.30\cdot 0.30 \\[2mm]
0.20
&\neq 0.09.
\end{align*}

So they are not independent.

\RepeatedQuestion{\uppgiftOneQuestion}
\textbf{b) Conditional probability}

We want:

\[
P(B\mid A^c).
\]

This means the probability of $B$ knowing that $A$ did not happen.

Use:

\begin{align*}
P(B\mid A^c)
&=
\frac{P(B\cap A^c)}{P(A^c)}.
\end{align*}

First:

\begin{align*}
P(A^c)
&=
1-P(A) \\[2mm]
&=
1-0.30 \\[2mm]
&=0.70.
\end{align*}

Then:

\begin{align*}
P(B\mid A^c)
&=
\frac{P(B\cap A^c)}{P(A^c)} \\[2mm]
&=
\frac{0.10}{0.70} \\[2mm]
&=
\frac{1}{7} \\[2mm]
&\approx 14.3\,\%.
\end{align*}

\section*{Uppgift 2}

\subsection*{Original question}
Låt följande funktion vara given.

\[
f(x)=
\begin{cases}
ax^2(2-x),
& \text{om } 0\leq x\leq 2,\\
0,
& \text{annars.}
\end{cases}
\]

\begin{enumerate}[label=\alph*)]
\item Bestäm talet $a$ sådant att $f$ är en täthetsfunktion.
\hfill (1,5 p)
\end{enumerate}

Givet att talet $a$ är sådant att $f$ är en täthetsfunktion,
låt $X$ vara en stokastisk variabel vars fördelning beskrivs av $f$.

\begin{enumerate}[label=\alph*),resume]
\item Bestäm fördelningsfunktionen för $X$.
\hfill (1,5 p)
\item Beräkna $P(X>1)$.
\hfill (1 p)
\item Beräkna väntevärdet och standardavvikelsen för $X$.
\hfill (2 p)
\end{enumerate}

\textit{Kommentar:}
Om du inte har löst a), så kan du istället använda
täthetsfunktionen
\[
f(x)=\frac{3x^2(2-x)}{4}
\]
för $0\leq x\leq 2$ för att lösa b)--d).

\subsection*{Compiled notes from Yacoub + Obid}

\subsubsection*{What we are doing}
The notes treat $f(x)$ as a continuous probability
density function.
For a continuous distribution, probability is represented
by the area under the curve.

So first we make the total area equal to $1$.
Then we integrate the PDF to get the CDF.
After that we use the CDF for $P(X>1)$.
Finally we find the expected value and the standard deviation.

\subsubsection*{Given information}
The given function is

\[
f(x)=
\begin{cases}
ax^2(2-x),
& 0\leq x\leq 2,\\
0,
& \text{otherwise.}
\end{cases}
\]

The notes also rewrite the inside as

\[
x^2(2-x)=2x^2-x^3.
\]

\subsubsection*{Step-by-step solution}

\RepeatedQuestion{\uppgiftTwoQuestion}
\textbf{a) Find $a$ so that $f$ is a PDF}

For a density function, the total area must be $1$:

\[
\int_{-\infty}^{\infty} f(x)\,dx=1.
\]

Since $f(x)$ is only nonzero from $0$ to $2$:

\begin{align*}
\int_0^2 ax^2(2-x)\,dx
&=1
\end{align*}

Expand:

\[
ax^2(2-x)=2ax^2-ax^3.
\]

Integrate:

\begin{align*}
\int_0^2 ax^2(2-x)\,dx
&=
\left[
\frac{2ax^3}{3}
-\frac{ax^4}{4}
\right]_0^2 \\[2mm]
&=1
\end{align*}

Plug in the bounds:

\begin{align*}
\left(
\frac{16a}{3}
-\frac{16a}{4}
\right)-0
&=1 \\[2mm]
\frac{64a-48a}{12}
&=1 \\[2mm]
\frac{16a}{12}
&=1 \\[2mm]
\frac{4a}{3}
&=1 \\[2mm]
a
&=\frac{3}{4}
\end{align*}

So the density becomes

\[
f(x)=\frac{3}{4}x^2(2-x),
\qquad 0\leq x\leq 2.
\]

\RepeatedQuestion{\uppgiftTwoQuestion}
\textbf{b) Find the distribution function}

The notes say that the cumulative distribution function
(CDF) is the probability that $X\leq x$.
It is found by integrating the PDF from the lower bound up
to $x$.

For $0\leq x\leq 2$:

\begin{align*}
F(x)
&=
\int_0^x
\frac{3}{4}t^2(2-t)\,dt \\[2mm]
&=
\frac{3}{4}
\left[
\frac{2t^3}{3}
-\frac{t^4}{4}
\right]_0^x
\end{align*}

Simplify:

\begin{align*}
F(x)
&=
\frac{x^3}{2}
-\frac{3x^4}{16},
\qquad 0\leq x\leq 2.
\end{align*}

So the full CDF is

\[
F(x)=
\begin{cases}
0,
& x<0,\\
\dfrac{x^3}{2}
-\dfrac{3x^4}{16},
& 0\leq x\leq 2,\\
1,
& x>2.
\end{cases}
\]

\RepeatedQuestion{\uppgiftTwoQuestion}
\textbf{c) Find $P(X>1)$}

The notes use the CDF from part b:

\begin{align*}
P(X>1)
&=1-P(X\leq 1) \\[2mm]
&=1-F(1)
\end{align*}

Plug in $x=1$:

\begin{align*}
F(1)
&=
\frac{1^3}{2}
-\frac{3(1)^4}{16} \\[2mm]
&=
\frac{8}{16}
-\frac{3}{16} \\[2mm]
&=
\frac{5}{16}
\end{align*}

Therefore:

\begin{align*}
P(X>1)
&=1-\frac{5}{16} \\[2mm]
&=\frac{11}{16} \\[2mm]
&=0.6875
\end{align*}

\RepeatedQuestion{\uppgiftTwoQuestion}
\textbf{d) Expected value and standard deviation}

The notes describe $E[X]$ as the theoretical average,
or the balance point of the distribution.

The formula is:

\[
E[X]=\int x f(x)\,dx.
\]

Use the density from part a:

\begin{align*}
E[X]
&=
\int_0^2
x\cdot \frac{3}{4}
\left(2x^2-x^3\right)\,dx \\[2mm]
&=
\int_0^2
\frac{3}{4}
\left(2x^3-x^4\right)\,dx
\end{align*}

Integrate:

\begin{align*}
E[X]
&=
\frac{3}{4}
\left[
\frac{x^4}{2}
-\frac{x^5}{5}
\right]_0^2 \\[2mm]
&=
\frac{3}{4}
\left(
\frac{16}{2}
-\frac{32}{5}
\right) \\[2mm]
&=
\frac{3}{4}(8-6.4) \\[2mm]
&=
\frac{3}{4}\cdot 1.6 \\[2mm]
&=1.2
\end{align*}

For the standard deviation, the notes first find
the variance:

\[
V[X]=E[X^2]-(E[X])^2.
\]

First find $E[X^2]$:

\begin{align*}
E[X^2]
&=
\int_0^2
x^2\cdot
\frac{3}{4}(2x^2-x^3)\,dx \\[2mm]
&=1.6
\end{align*}

\textbf{English clarification:}
The notes skip the middle integration for $E[X^2]$.
The same calculation is:

\begin{align*}
E[X^2]
&=
\frac{3}{4}
\int_0^2
(2x^4-x^5)\,dx \\[2mm]
&=
\frac{3}{4}
\left[
\frac{2x^5}{5}
-\frac{x^6}{6}
\right]_0^2 \\[2mm]
&=
\frac{3}{4}
\left(
\frac{64}{5}
-\frac{64}{6}
\right) \\[2mm]
&=1.6
\end{align*}

Now use the variance formula:

\begin{align*}
V[X]
&=
E[X^2]-(E[X])^2 \\[2mm]
&=
1.6-(1.2)^2 \\[2mm]
&=
1.6-1.44 \\[2mm]
&=0.16
\end{align*}

Standard deviation is the square root of the variance:

\begin{align*}
\sigma
&=
\sqrt{V[X]} \\[2mm]
&=
\sqrt{0.16} \\[2mm]
&=0.4
\end{align*}

\subsubsection*{Final answer}
\begin{enumerate}[label=\alph*)]
\item
\[
a=\frac{3}{4}.
\]

\item
\[
F(x)=
\begin{cases}
0,
& x<0,\\
\dfrac{x^3}{2}
-\dfrac{3x^4}{16},
& 0\leq x\leq 2,\\
1,
& x>2.
\end{cases}
\]

\item
\[
P(X>1)=\frac{11}{16}=0.6875.
\]

\item
\[
E[X]=1.2,
\qquad
\sigma=0.4.
\]
\end{enumerate}

\subsubsection*{What to remember}
For a PDF, the total area under the curve must be $1$.

For a CDF, integrate the PDF from the lower bound up to $x$.

Expected value is the balance point:

\[
E[X]=\int x f(x)\,dx.
\]

Variance and standard deviation measure the spread:

\[
V[X]=E[X^2]-(E[X])^2,
\qquad
\sigma=\sqrt{V[X]}.
\]

\section*{Uppgift 3}

\subsection*{Original question}
Låt $X \in \operatorname{Exp}(0,04)$ och
$Y \in \operatorname{Exp}(0,05)$ vara oberoende stokastiska variabler
och sätt
\[
Z=\min\{X,Y\}.
\]

\begin{enumerate}[label=\alph*)]
\item Bestäm fördelningen för $Z$.
\hfill (4 p)
\item Beräkna $P(Z\geq 10)$.
\hfill (1 p)
\item Bestäm väntevärdet för $Z$.
\hfill (1 p)
\end{enumerate}

\subsection*{Compiled notes from Yacoub + Obid}

\subsubsection*{What we are doing}
This problem is about the exponential distribution.
The notes describe it as a mathematical model for
how long we wait for an event to happen.

Here:

\[
X\sim \operatorname{Exp}(0.04)
\]

\[
Y\sim \operatorname{Exp}(0.05)
\]

The rates are usually represented by $\lambda$.
The new variable is

\[
Z=\min\{X,Y\}.
\]

So $Z$ means the time until the first one of the two
events happens.

\subsubsection*{Given information}
\begin{itemize}
\item Event $X$ happens at rate $0.04$ per unit of time.
\item Event $Y$ happens at rate $0.05$ per unit of time.
\item $X$ and $Y$ are independent.
\item $Z$ is the minimum of $X$ and $Y$.
\end{itemize}

The notes use the real-world picture:
imagine two light bulbs are running at the same time.
One burns out with rate $0.04$ and the other with rate $0.05$.
The minimum is how long until the first light bulb runs out.

\subsubsection*{Step-by-step solution}

\RepeatedQuestion{\uppgiftThreeQuestion}
\textbf{a) Determine the distribution for $Z$}

Because both events are running at the same time,
the rate at which at least one event happens is faster
than either rate alone.

For exponential distributions, the notes say that for
the first event, we add the individual rates together:

\begin{align*}
\lambda_Z
&=0.04+0.05 \\[2mm]
&=0.09
\end{align*}

Therefore the variable $Z$ follows an exponential distribution
with combined rate $0.09$:

\[
Z\sim \operatorname{Exp}(0.09).
\]

\RepeatedQuestion{\uppgiftThreeQuestion}
\textbf{b) Calculate $P(Z\geq 10)$}

The notes interpret this as:
what is the probability that we have to wait at least
$10$ units of time, or longer, before the first event happens?

For an exponential distribution, the notes use the formula
for waiting longer than $t$:

\[
P(Z\geq t)=e^{-\lambda t}.
\]

Plug in $\lambda=0.09$ and $t=10$:

\begin{align*}
P(Z\geq 10)
&=e^{-0.09\cdot 10} \\[2mm]
&=e^{-0.9} \\[2mm]
&\approx 0.4066
\end{align*}

So the probability of it happening after at least
$10$ units of time is about $40.66\,\%$.

\RepeatedQuestion{\uppgiftThreeQuestion}
\textbf{c) Determine the expected value for $Z$}

The notes describe the expected value as the mathematical
term for the average.

It asks:
on average, how long will it take for this event to happen?

For an exponential distribution, the formula for the expected
value is:

\[
\text{expected value}
=\frac{1}{\text{rate}}
=\frac{1}{\lambda}.
\]

Here $\lambda=0.09$, so:

\begin{align*}
E[Z]
&=\frac{1}{0.09} \\[2mm]
&=\frac{1}{9/100} \\[2mm]
&=\frac{100}{9} \\[2mm]
&\approx 11.11
\end{align*}

\subsubsection*{Final answer}
\begin{enumerate}[label=\alph*)]
\item
\[
Z\sim \operatorname{Exp}(0.09).
\]

\item
\[
P(Z\geq 10)=e^{-0.9}\approx 0.4066.
\]
That is about $40.66\,\%$.

\item
\[
E[Z]=\frac{1}{0.09}
=\frac{100}{9}
\approx 11.11.
\]
\end{enumerate}

\subsubsection*{What to remember}
For independent exponential waiting times,
the time until the first event is also exponential.

For the minimum, add the rates:

\[
\lambda_Z=\lambda_X+\lambda_Y.
\]

For waiting longer than $t$:

\[
P(Z\geq t)=e^{-\lambda t}.
\]

For the expected value:

\[
E[Z]=\frac{1}{\lambda}.
\]

\section*{Uppgift 4}

\subsection*{Original question}
\begin{enumerate}[label=\alph*)]
\item
I en förening genomförs en årlig pengainsamling, där medlemmar
ombeds att bidra med endera 50 kr eller 100 kr. Under de föregående
åren har $46\,\%$ valt att inte lämna något bidrag och andelen bidrag
om 50 kr har varit dubbelt så stor som andelen bidrag om 100 kr.
Låt $X$ vara lika med det bidraget från en slumpmässigt vald medlem
baserat på föregående år. Bestäm fördelningen för $X$ och beräkna
dess väntevärde samt standardavvikelse.
\hfill (2 p)

\item
Ett taxibolag prissätter taxiresor på följande sätt: varje resa har
en startavgift på 57 kr och sedan tillkommer 14 kr per km samt
9 kr per minut under resan. Enligt en undersökning är bolagets
innerstadsresor $X$ km långa och varar i $Y$ minuter, där
$X\in N(3,1)$ och $Y\in N(8,2)$. Beräkna sannolikheten för att en
slumpmässigt vald innerstadsresa kostar mellan 150 kr och 200 kr.
\hfill (4 p)
\end{enumerate}

\subsection*{Compiled notes from Yacoub + Obid}

\subsubsection*{What we are doing}
Part a is a discrete distribution problem.
The possible outcomes are:

\[
0\text{ kr},\qquad 50\text{ kr},\qquad 100\text{ kr}.
\]

The notes first find the probabilities for these outcomes.
Then they use the formulas for expected value, variance,
and standard deviation.

Part b creates a new normal variable for the taxi cost.
Then the notes standardize with

\[
Z=\frac{K-\mu}{\sigma}
\]

to find the probability between 150 kr and 200 kr.

\subsubsection*{Given information}
\textbf{Part a:}

\begin{itemize}
\item $46\,\%$ gave 0 kr.
\item The chance of 50 kr is twice as large as the chance of 100 kr.
\item $X$ is the contribution from one randomly selected member.
\end{itemize}

\textbf{Part b:}

\begin{itemize}
\item Start cost: 57 kr.
\item Cost per km: 14 kr.
\item Cost per minute: 9 kr.
\item $X\in N(3,1)$.
\item $Y\in N(8,2)$.
\end{itemize}

The notes use the first number as the mean and the second number
as the standard deviation:

\[
\mu_X=3,\quad \sigma_X=1,
\qquad
\mu_Y=8,\quad \sigma_Y=2.
\]

\subsubsection*{Step-by-step solution}

\RepeatedQuestion{\uppgiftFourQuestion}
\textbf{a) Distribution, expected value, and standard deviation}

Possible outcomes:

\[
X=0,\qquad X=50,\qquad X=100.
\]

From the question:

\[
P(X=0)=0.46.
\]

So the chance for people who paid something is:

\begin{align*}
1-0.46
&=0.54
\end{align*}

Let

\[
P(X=100)=p.
\]

The notes say that the chance of 50 kr is twice as large:

\[
P(X=50)=2p.
\]

Together they must be $0.54$:

\begin{align*}
p+2p
&=0.54 \\[2mm]
3p
&=0.54 \\[2mm]
p
&=\frac{0.54}{3} \\[2mm]
p
&=0.18
\end{align*}

Therefore:

\[
P(X=100)=0.18
\]

\[
P(X=50)=0.36
\]

\[
P(X=0)=0.46
\]

So the distribution is:

\[
\begin{array}{c|c}
x & P(X=x)\\
\hline
0 & 0.46\\
50 & 0.36\\
100 & 0.18
\end{array}
\]

Expected value means the average per person:

\begin{align*}
E[X]
&=
(0\cdot 0.46)
+(50\cdot 0.36)
+(100\cdot 0.18) \\[2mm]
&=
0+18+18 \\[2mm]
&=36
\end{align*}

So the expected value is:

\[
E[X]=36\text{ kr}.
\]

For standard deviation, first find $E[X^2]$:

\begin{align*}
E[X^2]
&=
(0^2\cdot 0.46)
+(50^2\cdot 0.36)
+(100^2\cdot 0.18) \\[2mm]
&=
0+900+1800 \\[2mm]
&=2700
\end{align*}

Variance:

\begin{align*}
\operatorname{Var}(X)
&=
E[X^2]-(E[X])^2 \\[2mm]
&=
2700-36^2 \\[2mm]
&=
2700-1296 \\[2mm]
&=1404
\end{align*}

Standard deviation:

\begin{align*}
\sigma
&=
\sqrt{\operatorname{Var}(X)} \\[2mm]
&=
\sqrt{1404} \\[2mm]
&\approx 37.47
\end{align*}

The notes mention that the standard deviation is bigger than
the mean, which tells us that the graph is spread out.

\RepeatedQuestion{\uppgiftFourQuestion}
\textbf{b) Taxi cost probability}

The taxi cost is written as:

\[
K=57+14X+9Y.
\]

The notes first find the mean:

\begin{align*}
\mu_K
&=
57+14\mu_X+9\mu_Y \\[2mm]
&=
57+14(3)+9(8) \\[2mm]
&=
57+42+72 \\[2mm]
&=171
\end{align*}

Now find the variance.
The notes say:
when combining independent variables, variances add,
constants disappear, and coefficients are squared.

\begin{align*}
\operatorname{Var}(K)
&=
14^2\operatorname{Var}(X)
+9^2\operatorname{Var}(Y) \\[2mm]
&=
14^2\cdot 1^2
+9^2\cdot 2^2 \\[2mm]
&=
196+324 \\[2mm]
&=520
\end{align*}

So:

\begin{align*}
\sigma_K
&=
\sqrt{520} \\[2mm]
&\approx 22.80
\end{align*}

Therefore:

\[
K\in N(171,22.80).
\]

We want:

\[
P(150\leq K\leq 200).
\]

Standardize:

\[
Z=\frac{K-\mu}{\sigma}.
\]

Lower value:

\begin{align*}
z_1
&=
\frac{150-171}{22.80} \\[2mm]
&\approx -0.92
\end{align*}

Upper value:

\begin{align*}
z_2
&=
\frac{200-171}{22.80} \\[2mm]
&\approx 1.27
\end{align*}

From the notes/table:

\begin{align*}
\Phi(z_1)
&=
\Phi(-0.92) \\[2mm]
&=
1-\Phi(0.92) \\[2mm]
&=
1-0.8212 \\[2mm]
&=0.1788
\end{align*}

Also:

\[
\Phi(z_2)=\Phi(1.27)=0.8980.
\]

So:

\begin{align*}
P(150\leq K\leq 200)
&=
\Phi(z_2)-\Phi(z_1) \\[2mm]
&=
0.8980-0.1788 \\[2mm]
&=0.7192
\end{align*}

\subsubsection*{Final answer}
\begin{enumerate}[label=\alph*)]
\item The distribution is:
\[
\begin{array}{c|c}
x & P(X=x)\\
\hline
0 & 0.46\\
50 & 0.36\\
100 & 0.18
\end{array}
\]

\[
E[X]=36\text{ kr}
\]

\[
\sigma\approx 37.47\text{ kr}.
\]

\item
\[
P(150\leq K\leq 200)
=0.7192.
\]

The probability for a taxi ride to cost between
150 kr and 200 kr is:

\[
71.92\,\%.
\]
\end{enumerate}

\subsubsection*{What to remember}
For a discrete expected value:

\[
E[X]=
\sum x\cdot P(X=x).
\]

For variance:

\[
\operatorname{Var}(X)
=E[X^2]-(E[X])^2.
\]

For a linear combination of independent normal variables:

\[
\operatorname{Var}(aX+bY)
=a^2\operatorname{Var}(X)
+b^2\operatorname{Var}(Y).
\]

To use the normal table:

\[
Z=\frac{K-\mu}{\sigma}.
\]

\section*{Uppgift 5}

\subsection*{Original question}
I ett slumpmässigt försök räknar en basketspelare hur många lyckade
straffkast hen kan göra i rad innan det första misslyckade
straffkastet. Låt $p$ vara sannolikheten för ett lyckat straffkast
och antag att alla straffkast är oberoende av varandra. Antalet
lyckade straffkast i försöket är då en stokastisk variabel $X$ med
följande sannolikhetsfunktion.

\[
p_X(k)=p^k(1-p),
\qquad
k=0,1,2,3,\ldots
\]

\begin{enumerate}[label=\alph*)]
\item
Låt $x_1,x_2,\ldots,x_n$ vara ett observerat stickprov från
fördelningen för $X$. Härled ML-skattningen
$p^*_{\mathrm{obs}}$ av $p$.
\hfill (5 p)

\item
Basketspelaren genomförde försöket vid 8 separata tillfällen med
följande resultat: $2,0,9,1,6,7,5,3$. Beräkna ML-skattningen av
$p$ utifrån detta resultat.
\hfill (1 p)
\end{enumerate}

\subsection*{Compiled notes from Yacoub + Obid}

\subsubsection*{What we are doing}
The notes use maximum likelihood.
We first write the likelihood function.
Then we take the log-likelihood because it is easier to maximize.
After that we differentiate, set the derivative equal to zero,
and solve for $p$.

\subsubsection*{Given information}
For one observation:

\[
p_X(k)=p^k(1-p),
\qquad
k=0,1,2,\ldots
\]

For the sample:

\[
x_1,x_2,\ldots,x_n
\]

The notes use:

\[
\sum x_i
\]

as the total number of successful shots in all attempts.

\subsubsection*{Step-by-step solution}

\RepeatedQuestion{\uppgiftFiveQuestion}
\textbf{a) Derive the ML estimate}

\textbf{Step 1: Likelihood function}

For each observation:

\[
p_X(x_i)=p^{x_i}(1-p).
\]

The likelihood is the product:

\begin{align*}
L(p)
&=
p^{x_1}(1-p)\,
p^{x_2}(1-p)\cdots
p^{x_n}(1-p) \\[2mm]
&=
p^{\sum x_i}(1-p)^n
\end{align*}

\textbf{Step 2: Log-likelihood}

The notes say maximizing $L(p)$ is harder because it has exponents.
Maximizing the log-likelihood is easier.
Because log is increasing, both have the same maximum.

\begin{align*}
\ell(p)
&=
\ln L(p) \\[2mm]
&=
\ln\!\left(
p^{\sum x_i}(1-p)^n
\right)
\end{align*}

Use

\[
\ln(ab)=\ln(a)+\ln(b)
\]

and

\[
\ln(k^x)=x\ln(k).
\]

Then:

\begin{align*}
\ell(p)
&=
\left(\sum x_i\right)\ln(p)
+n\ln(1-p)
\end{align*}

\textbf{Step 3: Differentiate}

Differentiate with respect to $p$:

\begin{align*}
\frac{d}{dp}\ell(p)
&=
\frac{\sum x_i}{p}
-\frac{n}{1-p}
\end{align*}

Set the derivative equal to zero:

\begin{align*}
0
&=
\frac{\sum x_i}{p}
-\frac{n}{1-p}
\end{align*}

So:

\begin{align*}
\frac{\sum x_i}{p}
&=
\frac{n}{1-p}
\end{align*}

\textbf{Step 4: Solve for $p$}

Cross multiply:

\begin{align*}
(1-p)\sum x_i
&=
np
\end{align*}

Expand:

\begin{align*}
\sum x_i
-p\sum x_i
&=
np
\end{align*}

Put all $p$ terms on one side:

\begin{align*}
\sum x_i
&=
np+p\sum x_i \\[2mm]
\sum x_i
&=
p\left(n+\sum x_i\right)
\end{align*}

Isolate $p$:

\begin{align*}
p
&=
\frac{\sum x_i}
{n+\sum x_i}
\end{align*}

So the ML estimate is:

\[
p^*_{\mathrm{obs}}
=
\frac{\sum x_i}
{n+\sum x_i}.
\]

The notes explain:

\[
\sum x_i
= \text{total successful shots}
\]

and

\[
n
= \text{number of failures}.
\]

So

\[
n+\sum x_i
= \text{total shots}.
\]

\RepeatedQuestion{\uppgiftFiveQuestion}
\textbf{b) Plug in the observed results}

The observed values are:

\[
2,0,9,1,6,7,5,3.
\]

There were 8 separate attempts, so:

\[
n=8.
\]

The notes mark this as the number of failures, since each attempt
ends with one missed shot.

Add the successful shots:

\begin{align*}
\sum x_i
&=
2+0+9+1+6+7+5+3 \\[2mm]
&=33
\end{align*}

Plug into the ML estimate:

\begin{align*}
p^*_{\mathrm{obs}}
&=
\frac{33}{8+33} \\[2mm]
&=
\frac{33}{41} \\[2mm]
&\approx 0.8049
\end{align*}

\subsubsection*{Final answer}
\begin{enumerate}[label=\alph*)]
\item
\[
p^*_{\mathrm{obs}}
=
\frac{\sum x_i}
{n+\sum x_i}.
\]

\item
\[
p^*_{\mathrm{obs}}
=
\frac{33}{41}
\approx 0.8049.
\]
\end{enumerate}

\subsubsection*{What to remember}
For maximum likelihood:

\begin{itemize}
\item multiply probabilities to get $L(p)$,
\item take $\ln(L(p))$ to simplify,
\item differentiate,
\item set the derivative equal to zero,
\item solve for the unknown parameter.
\end{itemize}

Here the estimate becomes:

\[
\hat p
=
\frac{\text{successful shots}}
{\text{total shots}}.
\]

\section*{Uppgift 6}

\subsection*{Original question}
Ett nytt parkeringshus har byggts i Grönköping med hopp om att locka
fler personer till butikerna i centrum. Kostnaden för parkeringshuset
ska på sikt täckas av intäkterna från parkeringsavgifter. Under
30 veckodagar noterades de dagliga intäkterna från parkeringsavgifter,
vilka anses utgöra ett observerat stickprov från en normalfördelning.
Stickprovets medelvärde och standardavvikelse är 1260 kr respektive
150 kr.

\begin{enumerate}[label=\alph*)]
\item
Bestäm ett $90\,\%$ konfidensintervall för den genomsnittliga dagliga
intäkten från parkeringsavgifter baserat på stickprovet ovan.
\hfill (3 p)
\end{enumerate}

Den konsult som var rådgivande när bygget av parkeringshuset
planerades förutsåg att den genomsnittliga dagliga intäkten från
parkeringsavgifter kommer att vara 1280 kr.

\begin{enumerate}[label=\alph*),resume]
\item
Ange lämpliga hypoteser $H_0$ och $H_1$ för att pröva konsultens
hypotes. Motivera ditt svar!
\hfill (1 p)

\item
Genomför hypotesprövning enligt b) med felrisk $10\,\%$ och baserat
på det observerade stickprovet ovan. Verkar konsultens hypotes rimlig
eller kan den förkastas?
\hfill (2 p)
\end{enumerate}

\subsection*{Compiled notes from Yacoub + Obid}

\subsubsection*{What we are doing}
For part a, the notes calculate a confidence interval.
This is a range of values that we are reasonably sure contains
the true average daily income.

For parts b and c, the notes use hypothesis testing.
The consultant says the average daily income will be 1280 kr.
We test whether this looks reasonable based on the sample.

\subsubsection*{Given information}
From the question:

\begin{itemize}
\item Sample size:
\[
n=30
\]
\item Sample mean:
\[
\bar{x}=1260
\]
\item Sample standard deviation:
\[
s=150
\]
\item Confidence level:
\[
90\,\%
\]
\item Error risk:
\[
\alpha=0.10
\]
\item Consultant's claim:
\[
\mu=1280
\]
\end{itemize}

\subsubsection*{Step-by-step solution}

\RepeatedQuestion{\uppgiftSixQuestion}
\textbf{a) 90\% confidence interval}

Since the population standard deviation is unknown and we use
the sample standard deviation $s$, the notes use the
$t$-distribution.

Degrees of freedom:

\begin{align*}
df
&=n-1 \\[2mm]
&=30-1 \\[2mm]
&=29
\end{align*}

For a $90\,\%$ confidence interval, the notes use:

\[
t_{0.05}\approx 1.699
\]

\textbf{Note from the notes:}
If the course uses the $z$-table for $n\geq 30$, then
$z_{0.05}\approx 1.645$ could be used.
The final handwritten interval uses the $t$ value.

The standard error is:

\begin{align*}
SE
&=
\frac{s}{\sqrt{n}} \\[2mm]
&=
\frac{150}{\sqrt{30}} \\[2mm]
&\approx
\frac{150}{5.477} \\[2mm]
&\approx 27.386
\end{align*}

Margin of error:

\begin{align*}
ME
&=
t\cdot \frac{s}{\sqrt{n}} \\[2mm]
&=
1.699\cdot
\frac{150}{\sqrt{30}} \\[2mm]
&\approx
1.699\cdot 27.386 \\[2mm]
&\approx 46.53
\end{align*}

Construct the interval:

\begin{align*}
\bar{x}\pm ME
&=
1260\pm 46.53
\end{align*}

Lower bound:

\begin{align*}
1260-46.53
&=1213.47
\end{align*}

Upper bound:

\begin{align*}
1260+46.53
&=1306.53
\end{align*}

So the confidence interval is:

\[
[1213.47,\;1306.53].
\]

\RepeatedQuestion{\uppgiftSixQuestion}
\textbf{b) Hypotheses}

The notes say there are two hypotheses:
either the person in the question is correct or wrong.

The consultant claims:

\[
\mu=1280.
\]

So use:

\[
H_0:\mu=1280
\]

\[
H_1:\mu\neq 1280
\]

This is two-sided, because we test whether the consultant's value is
different from 1280, not only bigger or smaller.

\RepeatedQuestion{\uppgiftSixQuestion}
\textbf{c) Hypothesis test with error risk 10\%}

The notes use the interval from part a:

\[
[1213.47,\;1306.53].
\]

Now check whether 1280 falls inside the interval.

\[
1280\in [1213.47,\;1306.53]
\]

Yes, it does.

Therefore the consultant's claim is reasonable based on this sample.
We do not reject $H_0$.

\subsubsection*{Final answer}
\begin{enumerate}[label=\alph*)]
\item
\[
90\,\%\text{ confidence interval}
=
[1213.47,\;1306.53]\text{ kr}.
\]

\item
\[
H_0:\mu=1280,
\qquad
H_1:\mu\neq 1280.
\]

\item
Since
\[
1280\in [1213.47,\;1306.53],
\]
we do not reject $H_0$.
The consultant's hypothesis seems reasonable.
\end{enumerate}

\subsubsection*{What to remember}
For a confidence interval with unknown population standard deviation:

\[
\bar{x}\pm t\frac{s}{\sqrt{n}}.
\]

For a two-sided hypothesis test at error risk $10\,\%$,
we can compare the claimed value with the $90\,\%$ confidence interval.
If the claimed value is inside the interval, we do not reject $H_0$.

\end{document}
